\rtmtight
\begin{tblr}{width=\textwidth,colspec={|Q[c,m,2.1cm]|X[l,m]|},hlines,vlines,hspan=minimal,rowsep=1.5pt,colsep=3pt,row{1}={bg=headgray,font=\bfseries}}
\SetCell{c,m}\hfil\logo\hfil & \textbf{POLITEKNIK NEGERI MALANG}\par\textbf{JURUSAN TEKNOLOGI INFORMASI}\par\textbf{PROGRAM STUDI: D4 TEKNIK INFORMATIKA} \\
\SetCell[c=2]{l} \textbf{RENCANA TUGAS MAHASISWA (RTM) DAN RUBRIK PENILAIAN} & \\
Nama Mata Kuliah & Proyek Sistem Informasi \\
Kode Mata Kuliah & RTI254009 \quad \textbf{sks / Jam perminggu:} 3 SKS/6 jam \quad \textbf{Semester:} 4 \\
Dosen Pengampu & Tim Pengajar Program Studi D4 TI \\
\end{tblr}
\assessmentblock{Dokumen Perencanaan Proyek SI}{Case Method dan penugasan terstruktur}{SCPMK0801-03601}{Mahasiswa menyusun dokumen perencanaan proyek SI lengkap meliputi WBS, jadwal, anggaran, dan risk register.}{Menganalisis kebutuhan proyek, menyusun WBS, menjadwalkan dengan Gantt chart, mengestimasi anggaran, mengidentifikasi risiko, dan mendokumentasikan dalam format standar.}{Dokumen perencanaan proyek, WBS, Gantt chart, risk register, dan laporan estimasi anggaran.}{Kelengkapan dan ketepatan WBS dan jadwal & 6 & WBS mendekomposisi semua deliverable dengan tepat, jadwal realistis, dan milestone teridentifikasi dengan jelas. & WBS dan jadwal tersedia tetapi dekomposisi atau milestone belum lengkap dan konsisten. & WBS tidak mencerminkan lingkup proyek atau jadwal tidak realistis. \\ \\
Kualitas risk register dan anggaran & 5 & Risiko diidentifikasi lengkap dengan probabilitas, dampak, dan mitigasi; anggaran terinci dan realistis. & Risiko dan anggaran tersedia tetapi mitigasi atau rincian anggaran masih kurang. & Tidak ada risk register yang memadai atau anggaran tidak dapat dipertanggungjawabkan. \\ \\
Kelengkapan dan kualitas dokumen & 4 & Dokumen lengkap, terformat standar, konsisten, dan dapat langsung digunakan sebagai baseline proyek. & Dokumen sebagian besar lengkap tetapi format atau konsistensi masih kurang. & Dokumen tidak lengkap atau tidak memenuhi standar minimal perencanaan proyek. \\}

\assessmentblock{Sprint Implementasi Sistem Informasi}{Project Based Learning}{SCPMK0802-03602, SCPMK0607-03604}{Mahasiswa mengimplementasikan SI berbasis web secara iteratif melalui sprint, mengintegrasikan back-end, front-end, dan database.}{Melaksanakan sprint planning, implementasi fitur, code review, integrasi, testing, dan sprint review sesuai metodologi Scrum.}{Kode sumber sistem informasi, hasil sprint review, dokumentasi teknis, dan laporan retrospective.}{Fungsionalitas sistem yang diimplementasikan & 18 & Fitur sprint berfungsi sesuai user story, integrasi back-end, front-end, dan database berhasil, dan pengujian unit lulus. & Fitur utama berfungsi, tetapi integrasi atau beberapa pengujian masih belum selesai. & Fitur tidak berfungsi atau integrasi sistem gagal. \\ \\
Kualitas kode, dokumentasi, dan integrasi & 13 & Kode bersih, mengikuti konvensi, terdokumentasi, dan integrasi antar lapisan arsitektur konsisten. & Kode berfungsi dan terdokumentasi sebagian, tetapi kualitas atau konsistensi integrasi masih terbatas. & Kode sulit dipahami, tidak terdokumentasi, atau integrasi tidak terstruktur. \\ \\
Proses agile, review, dan kolaborasi tim & 9 & Sprint review terdokumentasi, retrospective menghasilkan perbaikan konkret, dan kolaborasi tim terkoordinasi baik. & Sprint review dilakukan, tetapi retrospective atau dokumentasi kolaborasi masih kurang mendalam. & Tidak ada sprint review yang bermakna atau kolaborasi tim tidak terkoordinasi. \\}

\assessmentblock{UTS Milestone Review}{UTS berbasis demo dan presentasi}{SCPMK0803-03603}{Mahasiswa mendemonstrasikan kemajuan proyek SI dan mempertahankan keputusan teknis dan manajerial yang telah diambil.}{Mempersiapkan demo produk, laporan kemajuan, dan presentasi yang mencakup pencapaian sprint, risiko aktual, dan rencana ke depan.}{Demo sistem yang berjalan, laporan kemajuan proyek, dan bahan presentasi.}{Kelengkapan demo fitur Sprint 1-2 & 7 & Demo menunjukkan semua fitur Sprint 1-2 yang berfungsi, alur pengguna jelas, dan dapat dioperasikan oleh penguji. & Demo menunjukkan sebagian besar fitur, tetapi beberapa alur atau fungsionalitas masih belum stabil. & Demo tidak dapat dilaksanakan atau banyak fitur yang tidak berfungsi. \\ \\
Monitoring kemajuan dan pengelolaan risiko & 5 & Kemajuan terdokumentasi akurat, risiko aktual teridentifikasi, dan rencana mitigasi diperbarui secara realistis. & Kemajuan terdokumentasi tetapi risiko aktual atau rencana mitigasi belum cukup diperbarui. & Kemajuan tidak terdokumentasi atau risiko tidak dikelola secara proaktif. \\ \\
Presentasi dan argumentasi teknis & 3 & Keputusan teknis dan manajerial dijelaskan dengan alasan yang kuat dan mahasiswa mampu menjawab pertanyaan dengan tepat. & Presentasi cukup jelas, tetapi beberapa keputusan belum dapat diargumentasikan dengan baik. & Tidak dapat menjelaskan keputusan teknis atau manajerial secara memadai. \\}

\assessmentblock{UAS Demo Produk dan Laporan Akhir Proyek SI}{UAS berbasis demo dan defense}{SCPMK0607-03604, SCPMK0803-03603}{Mahasiswa mendemonstrasikan sistem informasi final yang deployable, menyerahkan laporan akhir proyek, dan mempertahankan hasil kerja tim.}{Melakukan finalisasi sistem, menyiapkan environment demo, menyusun laporan akhir proyek, dan mempresentasikan hasil secara komprehensif.}{Sistem SI yang deployable, laporan akhir proyek lengkap, dokumentasi teknis, dan bahan presentasi defense.}{Kelengkapan dan kualitas sistem akhir & 13 & Sistem SI dapat didemonstrasikan end-to-end, semua fitur utama berfungsi, dan telah melalui UAT dengan hasil terdokumentasi. & Sistem sebagian besar berfungsi, tetapi beberapa fitur atau hasil UAT masih kurang lengkap. & Sistem tidak dapat didemonstrasikan atau tidak melewati UAT minimal. \\ \\
Kualitas laporan dan dokumentasi proyek & 9 & Laporan akhir lengkap, mencakup perencanaan, implementasi, pengujian, dan refleksi proyek secara runtut dan terstandar. & Laporan tersedia tetapi beberapa bagian penting atau kelengkapan dokumentasi masih kurang. & Laporan tidak lengkap atau tidak mencerminkan perjalanan proyek secara akurat. \\ \\
Defense dan penguasaan teknis proyek & 8 & Mahasiswa mampu menjelaskan arsitektur, keputusan desain, tantangan, dan solusi proyek dengan meyakinkan dan akurat. & Defense cukup baik, tetapi beberapa aspek teknis atau manajerial belum dapat dijelaskan dengan mendalam. & Tidak dapat mempertahankan keputusan teknis atau manajerial proyek secara memadai. \\}

\begin{tblr}{width=\textwidth,colspec={|Q[l,m,3.2cm]|X[l,m]|},hlines,vlines,hspan=minimal,row{1-3}={bg=headgray},rowsep=1.5pt,colsep=3pt}
Jadwal Pelaksanaan: & Minggu 1--17 sesuai RPS. Setiap evaluasi memiliki RTM dan rubrik multi-kriteria. \\
Lain-Lain Yang Diperlukan: & LMS, repositori kode sumber, framework web, perangkat manajemen proyek, dan perangkat presentasi. \\
Pustaka: & Mengacu pustaka utama dan pendukung pada bagian RPS. \\
\end{tblr}
